\section{Evaluation}
\label{sec:eval}

We evaluate the deployed system on three corpora differing in origin, language
and scale, and we separate two questions: whether the management layer keeps
its collections in a usable state, and what the answers built on those
collections are worth. Sections~\ref{sec:e1}--\ref{sec:e3} cover the
management layer: acquisition and parsing, organization and three-source
consistency, and the operational-experience store. Sections
\ref{sec:e4}--\ref{sec:e8} cover the query path: retrieval quality
(Section~\ref{sec:e4}), answer quality per question and abstention
(Section~\ref{sec:e5}, the paper's end-to-end evidence), verifier failure
(Section~\ref{sec:e6}), non-functional behavior (Section~\ref{sec:e7}) and
the gate and search-scope ablations (Section~\ref{sec:e8}). Numbers are
reported against the corpus that produced them, never merged; the campaign
inventory and per-question tool-layer results are supplementary material.

\subsection{Experimental setup}
\label{sec:eval-setup}

\textbf{Corpora.} (i)~\emph{Deployed industrial corpus}: five knowledge bases
in a live deployment hold 42 documents, one containing twenty near-duplicate
equipment-telemetry reports sharing prefixes, vocabulary and structure---a
deliberate stressor for near-duplicate detection and confuser
density~\cite{manku-near-dup-2007}; sixteen gold questions and four
out-of-corpus probes ran against it. (ii)~\emph{arXiv corpus}: 100 submitted
papers spanning about 25 domains, 79 retrievable at evaluation time and routed
by content into five domain libraries holding 322 stored parts; ten gold
questions (BQ01--BQ10) target papers referenced by those papers, keyed by arXiv
identifier, and the live environment held 19 libraries. (iii)~\emph{Cross-lingual
ingestion corpus}: 400 Wikipedia pages per language (English, Chinese,
Japanese), 40 extractive cloze queries each with verbatim in-page answers;
ingestion only, no retrieval-quality numbers reported.

\textbf{Questions and gold annotations.} The sixteen industrial questions were
authored, not sampled from logs; each identifies one or two gold documents by
a unique filename fragment and one to three fact markers. Q16 includes a token
also present in its gold filename, so a filename-ranking retriever could match
it without reading content. We retain this limitation. The arXiv answer grader
checks gold-ID presence in the answer or trace plus a word-boundary keyword
coverage rule; it does not assess semantic support. Citation markers are
counted separately by a deterministic pattern and do not assess source
resolution or claim-to-source entailment.

\textbf{Metric families.} Retrieval uses Hit@5, Recall@5 and MRR@5, following
passage-ranking conventions~\cite{craswell-msmarco-2021};
nDCG@5~\cite{jarvelin-cumulated-gain-2002} was computed only inside one
exploratory arm we do not report. The end-to-end "gold-hit" measure checks
whether the answer or recorded retrieval trace contains the gold identifier;
keyword coverage is a deterministic lexical proxy. Citation counts record
whether the answer contains a part, section or table marker under the stated
pattern. These checks support traceability and marker-presence claims, not
semantic faithfulness or claim-to-source support, for which model-based
measures such as RAGAS~\cite{es-ragas-2024} or ARES~\cite{saad-falcon-ares-2024}
are designed. The third metric family covers wall-clock time, agent turns,
tool calls, cost and whether refusals were appropriate.

\textbf{Controls.} Deterministic lanes were re-run and compared field by field,
latency excluded; model-driven lanes were repeated at fixed model version and
reported as means with observed spread; the fixed retry policy (one retry for
an empty zero-turn response or transient gateway timeout) fired twice,
succeeding both times. Before measurement we repaired two corpus-integrity
faults via the platform's own operations (Section~\ref{sec:e7}). The tool-layer
retrieval lanes, grading rule, audits, duplicate adjudication and cross-lingual
ingestion pipeline are deterministic; agent lanes and text-derived results are
not, and carry their repetition spread.

\subsection{E1: Acquisition and parsing}
\label{sec:e1}

The 100 submitted papers were parsed and content-routed into five domain
libraries, producing 322 stored parts after the 30{,}000-character split. The
submission funnel---91 in the ingest state file, 79 retrievable at evaluation
time, 21 never retrievable---is in the campaign inventory; reported failure
modes are formula-dense scans and one tool-layer crash during batch parsing,
restarted from a checkpoint. Neither attempts against failed papers nor
per-paper parse status is recorded, so reconciliation rests on surviving
artifacts~\cite{missier-prov-2013}. The audit: nine per-paper ingest checks
100/100; no true duplicate among 35 near-duplicate pairs; 165 descriptions
upgraded; content and storage audits 322/322 and 33/33; 15/15 regression
probes; a 4{,}528-chunk dense replica index (values script-asserted from saved
artifacts).

The cross-lingual benchmark ran twice over 400 pages per language and
reproduced round~1 bit-for-bit, with ingestion success 1.0; storage
completeness, self-retrieval Hit@1 and the membership figures are in
Table~S2. Its English membership accuracy of 0.875 is an organization
defect---the page stored in the wrong library roughly one time in eight---not
a parsing defect. Membership is scored against the benchmark's source
category labels rather than against any naming convention derived by the
ingest agent, so it is the one placement-correctness measurement in this
paper that does not depend on filename cues.

A 50-paper, 26-domain campaign parsed and verified 50/50 papers in
1{,}115.5\,s (8.3--48.4\,s per paper), produced 202 tags, batch-indexed 50
documents and reported \texttt{all\_probes\_ok}; 41 of 50 carried
\texttt{tag\_ok:false} although \texttt{verify\_ok} passed 50/50, so
tag-quality failures never blocked ingestion. Its 30{,}000-character tier
turned 50 papers into 165 documents (nine whole, 156 parts), passed storage
50/50 and the description-and-tag check 165/165, stored 3{,}798{,}446
characters and parsed in 2{,}672\,s; one stage re-ran on inbox ordering, and
more than 800 extracted images were not uploaded. Descriptions and tags are
checked against content before storage~\cite{park-metadata-quality-2009}.


\subsection{E2: Organization, consistency, and management operations}
\label{sec:e2}


Organizational correctness is three-source consistency: disk files, the tree
index and per-library metadata must agree on document existence, location and
indexed content. Under that definition we found one real defect in a
disposable test library: after a rename the vector index kept the old name's
chunks---count 18 where 9 were expected, a duplicate search matching the
document against its own former path at 0.9738 similarity, a vector probe
returning every chunk twice at identical scores. Only an explicit force
re-index removed them, restoring the count to 9. The failure is invisible from
the metadata alone.

The full reorganization campaign ran over the 322-document arXiv corpus and
passed every check in the campaign audit; its 35 near-duplicate candidate
groups (similarities 0.90--0.975), each adjacent parts of one paper, held zero
true duplicates, and no document was deleted. The spot check sampled 33 of the
322 documents. Sub-library splitting and hierarchy creation were executed
later on a disposable three-level fixture (2026-10-13): creation with a parent
reference, a parent$\to$sub-library document move, and listing, reading and
judging through sub-library addresses all returned correct results, and a
parent catalog lists its sub-libraries as container entries carrying the
sub-library's own description---the one-call view ladder navigation reads.

Growing from 50 to 100 papers (322 stored documents at the final audit)
cost no retrieval
quality: dense top-1 hits went from 48/50 to 49/50 and MRR from 0.98 to 1.00,
a 2 percentage point gain recorded as ``held (within $\pm$5\,pp)''. In the 16
lifecycle scenarios, run against two disposable test libraries, create, move,
delete and forced re-index each behaved as specified
(Table~S5). Two failure modes: the task-status tool reports
\texttt{running} after \texttt{finished\_at} is written, as in the
reorganization graph build, and the update check timed out after 30\,s
(\texttt{ragctl version -{}-json timed out after 30s}, success false) while
every other probe passed; the project-status tool silently ignores its version
flag. Merge and split were never exercised under controlled conditions and we
claim nothing about them; the ingest tools were covered in E1.

The live tag registry held 316 tags; the tag-quality check passed 165 documents
with zero failures in 47.6\,s. Orphan tags accumulate (test tags, section
titles); the registry offers a dry-run detector and a cleanup, and in one case
test-library tags leaked into the production registry---the failure mode
vocabulary-based organization is prone to~\cite{golder-tagging-2006}. The
document graph is navigation, never evidence: writing one document's four tag
edges moved the edge count from 4{,}675 to 4{,}679, against 658 nodes, 333
documents and 316 tags; the 165-document corpus graph has 335 nodes and
1{,}671 edges.
